% Historical retained-population insert, superseded for the main manuscript by
% evaluation_delivery/evaluation.tex and the fresh V9 domain ablation.
% Add to Section 5.1 after baseline definitions.
To assess the contribution of batching to hardware selection, we compare the
serial reference, joint exploration without cross-chunk batching, and full
RoboDSE. The no-batching arm retains serial execution, the $k=1$ pipeline,
and unbatched $k=2,3$ organizations. We report both selection over the frozen
retained candidates and a matched-support check requiring all five pipeline
organizations at each retained hardware, core-split, and relative-period-factor
combination. This check controls observed organization coverage but does not
recover mappings omitted by the original incumbent-based pruning.

% Add to Section 5.3 after the resource-frontier overview.
\paragraph{Contribution of batching to design selection.}
Over the frozen retained candidates, joint exploration without batching yields
the same core-count/bandwidth frontiers as the serial reference in all 22 boxes.
Full RoboDSE produces five strictly improved frontiers and three extensions
relative to the no-batching arm. Under matched organization support, the five
strict improvements remain, while the other 17 boxes have identical frontiers.
These classification counts also hold when conditional NoRefresh configurations
are excluded and when total SRAM is included as a third resource dimension.
Thus, the five strict improvements provide evidence for the contribution of
batching within the matched evaluated space. The three extensions observed in
the full retained population are not reproduced by this restricted-support
check and are not presented as matched-support batching attribution. Identical
resource frontiers do not imply identical execution timings; no exhaustive
unpruned optimum or task-success improvement is established by this ablation.
